\documentclass[12pt]{article}

\input{../../../_assets/preambulo.tex}
\input{../../../_assets/portada.tex}

\usepackage[all]{xy}
\usepackage{ stmaryrd }

\begin{document} 
    \portada[%
        titulo=Modelos Matemáticos II,
        subtitulo=,
        autor=Jesús Muñoz Velasco,
        año=Curso 2025-2026]
        
    \thispagestyle{empty}

    % os recuerdo que tendremos la primera prueba de clase el próximo miércoles
    % 13 de Mayo, en horario de clase. Durante la semana os envío un documento
    % con una convocatoria y os detallo instrucciones para el día de la prueba.

    % A modo de resumen, los contenidos tratados son los siguientes:

    % - Nociones sobre optimización en dimensión finita y programación lineal.

    % - Optimización de funcionales: el método directo; derivadas direccionales.

    % - Condiciones necesarias: ecuación de Euler-Lagrange (casos de funciones escalares de una y varias variables, y caso de curvas).

    % - Condiciones suficientes: ecuación de Jacobi.

    % - Nociones de convexidad/concavidad.

    % - Lagrangianos y principio de mínima acción.

    % - Optimización con ligaduras integrales.

    % - Problemas de valores en la frontera. Teorema de Fredholm.

    % - Problemas de Sturm-Liouville.

    % - Problema de Cauchy para la ecuación de ondas.

    % - Series de Fourier.

    % - Problema mixto para la ecuación de ondas.

    % Para cualquier consulta al respecto poneos en contacto conmigo. Un saludo,

    % Juan
    
    \section{Nociones sobre optimización en dimensión finita y programación lineal}

    \subsection{Weierstrass}
    Si $D\subset\mathbb R^n$ es compacto y $f:D\to\mathbb R$ es continua,
    entonces $f$ alcanza máximo y mínimo global.

    \subsection{Puntos críticos}
    Si $D$ es abierto y $f\in C^1(D)$ alcanza un extremo local en $x_0\in D$,
    entonces
    \[
    \nabla f(x_0)=0.
    \]

    \subsection{Multiplicadores de Lagrange en dimensión finita}
    Para optimizar $f(x)$ sujeto a $g(x)=0$, se resuelve
    \[
    \nabla f(x)=\lambda \nabla g(x),\qquad g(x)=0.
    \]

    \subsection{Programación lineal}
    En programación lineal, si el óptimo existe y el conjunto admisible es un poliedro,
    se alcanza en un vértice. Si el conjunto no está acotado, hay que comprobar si el
    máximo/mínimo existe.

    \section{Optimización de funcionales: el método directo; derivadas direccionales}

    \subsection{Método Directo}
    \begin{enumerate}
        \item Ver si el funcional está acotado
        \item En caso de estar acotado hallar la mejor cota posible y demostrar que se alcanza o que hay una sucesión que tiende a dicha cota en $D$ (si el límite no esta en $D$ entonces no se alcanza dicha cota).
    \end{enumerate}

    \subsection{Derivadas direccionales (Derivada de Gateaux)}
    \begin{gather*}
        \delta_\varphi \cc{F}[y] = [d_\veps \cc{F}[y+\veps \varphi]]_{|_{\veps = 0}}
    \end{gather*}
    \begin{enumerate}
        \item Perturbar $\varphi$ en una dirección admisible (genérica).
        \item Derivar con respecto a $\veps$.
        \item Evaluar en $\veps=0$.
    \end{enumerate}
    Si $y$ es extremo de $\mathcal F$, entonces
    \begin{align*}
        \delta_\varphi \mathcal F[y]=0
    \end{align*}
    para toda dirección admisible $\varphi$

    \section{Condiciones necesarias: ecuación de Euler-Lagrange (casos de funciones escalares de una y varias variables, y caso de curvas)}

    % D=[x_0, x_1]

    \subsection{Una variable}
    Consideramos el funcional dado por 
    \begin{gather*}
        \cc{F}[y] = \int_{x_0}^{x_1} F(x,y,p) dx
    \end{gather*}
    donde $y:[x_0, x_1]\to \bb{R}$. La Ecuación de Euler-Lagrange queda
    \begin{gather*}
        F_y -\frac{d}{dx}(F_p) = 0
    \end{gather*}
    Cuando los extremos no son fijos, a la ecuación de Euler-Lagrange hay que añadirle la condición
    \begin{gather*}
        F_p(x_0,y(x_0), y'(x_0)) = 0 = F_p(x_1,y(x_1), y'(x_1))
    \end{gather*}
    (si solo hubiera un extremo libre se aplicaría esta condición solo a dicho extremo).
    Cuando la condición es de tipo periódica ($y(x_0)=y(x_1)$) añadimos la condición
    \begin{gather*}
        F_p(x_0,y(x_0), y'(x_0)) = F_p(x_1,y(x_1), y'(x_1))
    \end{gather*}

    \subsubsection{Funcionales dependientes de $y''$}
    Si se tiene que el funcional está definido como 
    \begin{gather*}
        \mathcal F[y]=\int_{x_0}^{x_1}F(x,y,p,q)\,dx,
    \end{gather*}
    donde $p=y'$, $q=y''$ entonces la ecuación de Euler-Lagrange es
    \begin{gather*}
        F_y-\frac{d}{dx}F_{p}+\frac{d^2}{dx^2}F_{q}=0.
    \end{gather*}

    \subsection{Varias variables}
    Consideramos el funcional dado por 
    \begin{gather*}
        \cc{F}[y] = \int_\Omega F(x,y,\nabla y) dx = \int_\Omega F(x_1,...,x_d,y,p_1,...,p_d) dx_1...dx_d
    \end{gather*}
    donde $y: \Omega \subset \bb{R}^d \to \bb{R}$. La Ecuación de Euler-Lagrange queda
    \begin{gather*}
        F_y - div_{(x_1,...,x_d)}(F_p) = 0
    \end{gather*}
    donde
    \begin{gather*}
        div_x(F_p) = \sum_{i=1}^d \frac{\partial F_{p_i}}{\partial x_i} =
        \frac{\partial F_{p_1}}{\partial x_1} + ... + \frac{\partial F_{p_d}}{\partial x_d}
    \end{gather*}

    \subsection{Curvas}
    Consideramos el funcional dado por 
    \begin{gather*}
        \cc{F}[y] = \int_{x_0}^{x_1} F(x,y,p) dx = \int_{x_0}^{x_1} F(x, y_1,...,y_d, p_1,...,p_d)dx
    \end{gather*}
    donde $y:[x_0, x_1]\to \bb{R}^d$. La Ecuación de Euler-Lagrange es un sistema de $d$ ecuaciones dado por
    \begin{gather*}
        F_{y_i} - \frac{d}{dx}(F_{p_i})=0\ \ \ i\in\{1,...,d\}
    \end{gather*}

    \subsection{Identidad de Beltrami}
    Si $F$ no depende explícitamente de $x$, entonces
    \[
    F-y'F_p=C.
    \]

    \section{Condiciones suficientes: Ecuación de Jacobi}
    Se haría la segunda derivada de Gateaux. Por simplicidad solo estudiamos esta ecuación cuando se tienen extremos fijos en un intervalo $[x_0, x_1]$. El procedimiento para comprobar si un extremal $y$ es mínimo sería el siguiente:
    \begin{enumerate}
        \item Calculamos $P$ y $Q$ dados por 
        \begin{align*}
            P &= F_{pp}\\
            Q &= F_{yy}-\frac{d}{dx}(F_{yp})
        \end{align*}
        \item Comprobamos la \textbf{condición necesaria de Legendre} dada por 
        \begin{gather*}
            \begin{array}{l l l}
                \text{Si } P\geq 0 & \Rightarrow & \text{ Se cumple la condición}\\
                \text{Si } P > 0 & \Rightarrow & \text{ Se cumple la condición estricta}
            \end{array}
        \end{gather*}
        Si no se cumple ninguna de estas condiciones no podemos concluir nada con la condición de Jacobi.
        En caso contrario continuamos con el procedimiento.
        \item Resolvemos la \textbf{ecuación de Jacobi} con incógnita $\varphi \in C^2([x_0, x_1])$ dada por
        \begin{gather*}
            -\frac{d}{dx}(P\varphi') + Q\varphi = 0
        \end{gather*}

        \item 
        
        \begin{definicion}[Punto conjugado]
            Un punto $\bar x\in ]x_0,x_1]$ se dice \textbf{conjugado} de $x_0$ si existe una solución no trivial
            $\varphi$ de la ecuación de Jacobi tal que
            \begin{gather*}
                \varphi(x_0)=0,\qquad \varphi(\bar x)=0.
            \end{gather*}
            \endsquare
        \end{definicion}

        Si no existen puntos conjugados en $]x_0,x_1]$ y se cumple Legendre estricta, el extremal es mínimo local estricto. Esto se comprobará imponiendo
        \begin{gather*}
            \varphi(x_0) = 0
        \end{gather*}
        y tendremos que demostrar que o bien $\varphi(x)\equiv 0$ o bien $\nexists \bar{x}\in]x_0, x_1]$ con
        \begin{gather*}
            \varphi(\bar{x})=0
        \end{gather*}
    \end{enumerate} 

    \section{Nociones de convexidad/concavidad}

    \begin{definicion}[Subconjunto convexo]
        Tomamos un espacio vectorial real $X$, diremos que un subconjunto $U\subset X$ es convexo si para cualesquiera $u,v\in U$ se tiene que 
        \begin{gather*}
            \theta u + (1-\theta) v \in U \ \ \ \forall \theta \in[0,1]
        \end{gather*}
    \end{definicion}

    Se ha estudiado ya que cualquier conjunto de la forma
    \begin{gather}\label{def-dominio}
        \cc{D}= \{y\in C^1([x_0, x_1],\bb{R}^N) : y(x_0)=y_0,\ y(x_1)=y_1\},\ \ \ y_0,y_1\in \bb{R}^N \text{ fijos}
    \end{gather}
    es convexo.

    \begin{definicion}[Funcional convexo]
        Diremos que un funcional $\cc{F} : \cc{D} \to \bb{R}$ es convexo si para cualesquiera $u,v\in D$ se tiene que
        \begin{gather*}
            \cc{F}[\theta u + (1-\theta) v] \leq \theta\cc{F}[u] + (1-\theta) \cc{F}[v] \ \ \ \forall \theta \in [0,1]
        \end{gather*}
    \end{definicion}

    \begin{lema}
        Para cada $x\in I$ definimos $F_x:\mathbb{R}^N\times\mathbb{R}^N\to \mathbb{R}$ dada por:
        \begin{equation*}
            F_x(y,p) = F(x,y,p), \qquad (y,p)\in \mathbb{R}^N\times\mathbb{R}^N
        \end{equation*}
        Si para cada $x\in I$ tenemos que esta aplicación es convexa, entonces el funcional $\cc{F}:I\times\mathbb{R}^N\times \mathbb{R}^N\to \mathbb{R}$ dado por:
        \begin{equation*}
            \cc{F}(y) = \int_{x_0}^{x_1} F(x,y(x),y'(x))~dx 
        \end{equation*}
        es convexo en cualquier subconjunto convexo de $C^1(I, \mathbb{R}^N)$.
    \end{lema}

    \begin{observacion}
        Para ver que $F_x(y,p)$ es efectivamente convexo, lo que haremos será calcular
        \begin{gather*}
            Hess F_x(y,p) = 
            \begin{pmatrix}
                \dfrac{\partial^2 F_x}{\partial y^2} & \dfrac{\partial^2 F_x}{\partial y \partial p}\\
                \dfrac{\partial^2 F_x}{\partial p \partial y} & \dfrac{\partial^2 F_x}{\partial p^2}
            \end{pmatrix} (y,p)
        \end{gather*}
        Ahora haremos una distinción de casos:
        \begin{itemize}
            \item Si $Hess F_x(y,p)$ es semidefinida positiva, entonces el funcional $\cc{F}$ es convexo.
            \item Si $Hess F_x(y,p)$ es definida positiva, entonces el funcional $\cc{F}$ es estrictamente convexo.
            \item Si $Hess F_x(y,p)$ es semidefinida negativa, entonces el funcional $\cc{F}$ es cóncavo.
            \item Si $Hess F_x(y,p)$ es definida negativa, entonces el funcional $\cc{F}$ es estrictamente cóncavo.
        \end{itemize}
        Para comprobar que la matriz es definida positiva (equivalentemente para negativa) bastará comprobar los menores principales líderes. \\

        Para comprobar que la matriz es semidefinida positiva (equivalentemente para negativa) no basta siempre con mirar solo los menores principales líderes; hay que comprobar los menores principales.
        En dimensión 2, para
        \[
        H=\begin{pmatrix}a&b\\ b&c\end{pmatrix},
        \]
        semidefinida positiva equivale a
        \[
        a\geq 0,\qquad c\geq 0,\qquad |H|\geq 0.
        \]
        
        Las consecuencias sobre concavidad solo funcionarán en subconjuntos convexos de $C^1(I, \bb{R}^N)$, como indica el lema. Como la mayoría serán de la forma (\ref{def-dominio})
    \end{observacion}

    \begin{itemize}
        \item Si $\cc{F}$ es estrictamente convexo, a lo sumo habrá un mínimo.
        \item Si $\cc F$ es convexo y $y$ es punto crítico, entonces $y$ es mínimo global.
        \item Si $\cc F$ es estrictamente convexo, dicho mínimo global es único
    \end{itemize}

    \section{Lagrangianos y principio de mínima acción.}

    \section{Optimización con ligaduras integrales.}

    Tenemos un problema consistente en calcular los extremales de un funcional $\cc{F}:\cc{D} \to \bb{R}$ dado por 
    \begin{gather*}
        \cc{F}[y] = \int_{x_0}^{x_1} F(x,y,p) dx
    \end{gather*}
    sobre un dominio de la forma 
    \begin{gather*}
        \cc{D} = \left\{ y\in C^1([x_0, x_1]) y(x_0)=y_0,\ y(x_1)=y_1,\ \int_{x_0}^{x_1} \Phi(x,y,p) dx = \varphi \right\}
    \end{gather*}
    En general, las condiciones del dominio pueden incluir otra condición integral y las condiciones de tipo Dirichlet pueden ser de otro tipo. \\

    \begin{enumerate}
        \item Introducimos el funcional corregido como 
        \begin{gather*}
            F^*(x,y,p) = F(x,y,p) - \lambda \Phi(x,y,p)
        \end{gather*}

        En caso de que haya varias condiciones integrales (supongamos que hay $n$), el funcional extendido será
        \begin{gather*}
            F^*(x,y,p) = F(x,y,p) - \lambda_1 \Phi_1(x,y,p) - ... - \lambda_n \Phi_n(x,y,p)
        \end{gather*}
        \item Le aplicamos Euler-Lagrange al funcional corregido y resolvemos la ecuación diferencial en función de $\lambda$ (o de los múltiples $\lambda$ en caso de haberlos). 
        \item Imponemos las condiciones de contorno para obtener una solución, que seguirá en función de $\lambda$.
        \item Aplicamos las condiciones integrales a la solución obtenida, llegando al valor de $\lambda$ (es posible que durante el proceso lleguemos a un problema de Sturm-Liouville).
        \item Sustituimos el valor de $\lambda$ en la solución $y$ y esta será la solución del problema estudiado.
    \end{enumerate}
 
    \section{Problemas de valores en la frontera. Teorema de Fredholm}
    Consideramos un problema de contorno de la forma 
    \begin{gather*}
        \left\{
        \begin{array}{l}
            (p(x)y'(x))' + q(x)y(x) = r(x)\\
            y(x_0) = y(x_1) = 0
        \end{array}
        \right.
    \end{gather*}   
    donde $0<p\in C^1([x_0, x_1])$, $q,r\in C([x_0, x_1])$. 
    Nuestro objetivo es estudiar cuántas soluciones puede tener el problema
    \begin{enumerate}
        \item Consideramos el problema homogéneo eliminando $r(x)$ de la siguiente forma:
        \begin{gather*}
            \left\{
            \begin{array}{l}
                (p(x)y'(x))' + q(x)y(x) = 0\\
                y(x_0) = y(x_1) = 0
            \end{array}
            \right.
        \end{gather*}  
        \item Resolvemos el problema de contorno.
        \begin{itemize}
            \item Si solo existe una solución y es trivial, entonces el problema original tiene una única solución y habremos terminado el estudio.
            \item En otro caso el problema completo tendrá solución si y solo si todo $\psi$ solución del problema homogéneo es ortogonal a $r$, es decir,
            \begin{gather*}
                \langle r, \psi \rangle = \int_{x_0}^{x_1} r(x)\psi(x)dx = 0
            \end{gather*}
            En este caso el problema completo tendrá infinitas soluciones. En caso de que no sean ortogonales el problema original no tendrá ninguna solución.
        \end{itemize}
    \end{enumerate}
    \begin{observacion}
        Una forma fácil de ver que el problema homogéneo solo admite la solución trivial es construir el funcional asociado
        \begin{gather*}
            \cc{F}[y] = \int_{x_0}^{x_1} p \frac{y'^2}{2} - q \frac{y^2}{2} dx
        \end{gather*}
        que tendrá como ecuación de Euler-Lagrange la ecuación diferencial que tenía el problema homogéneo. Si conseguimos demostrar que el funcional es estrictamente convexo y que la solución trivial es solución habríamos probado que es la única posible.
    \end{observacion}

    \section{Problema mixto para la ecuación de ondas.}

    Tenemos el problema:
    \begin{equation*}
        \left\{\begin{array}{ll}
                u_{tt} - c^2 u_{xx} = 0 , & \text{en\ } \left[0,\infty\right[\times [0,L] \\
                    u(0,x) = \varphi(x) , &\text{en\ } [0,L] \\
                    u_t(0,x) = \psi(x) , &\text{en\ } [0,L] \\
                    u(t,0) = 0 = u(t,L) , & \text{en\ } \left[0,\infty\right[
        \end{array}\right.
    \end{equation*}

    La estrategia de resolución es la siguiente:
    \begin{enumerate}
        \item Primero resolvemos por variables separadas $u(t,x) = T(t) X(x)$, y expresamos
        $u_{tt} - c^2 u_{xx} = 0$ con la nueva descomposición como dos cocientes iguales, 
        uno dependiendo solo de $t$ y otro solo de $x$, por lo que ambos deben ser constantes,
        pongamos a $\lm \in \mathbb{R}$. 
        \item Distinguimos casos comparando $\lm$ con $0$:
        \begin{itemize}
            \item Si $\lm < 0$ obtenemos combinación de exponenciales, imponemos condiciones de contorno y vemos qué solución queda.
            \item Si $\lm = 0$ resolvemos $X_0'' = 0$, luego $X_0(x) = ax+b$, para $a,b\in \mathbb{R}$ y obtenemos imponiendo las condiciones de contorno la solución correspondiente.
            \item Si $\lm > 0$ obtenemos combinación de senos y cosenos de la misma frecuencia, cosa que se desarrollará en el cálculo a continuación.
        \end{itemize}
        \item Se obtendrán, para cada $n \in \mathbb{N}$, soluciones $T_n(t)$ y $X_n(x)$, y a partir de estas dos, $u_n(t,x)$. Así, la solución más general posible (puesto que la ecuación es lineal) es:
        \begin{equation*}
            u(t,x) = \sum_{n=1}^{\infty} u_n(t,x)
        \end{equation*} 
        \item Finalmente se imponen las condiciones de contorno $$u(0,x) = \varphi(x) \quad u_t(0,x) = \psi(x))$$
        para ciertas $\varphi, \psi$ dadas, y se identifican los sumandos no nulos, obteniendo así la solución final.
    \end{enumerate}

    Veamos un ejemplo detallado siguiendo este esquema: \\

    Por variables separadas, $u(t,x) = T(t) X(x)$. Sustituimos la condición primera del problema, 
    $$u_{tt} - c^2 u_{xx} = 0 \iff T''(t) X(x) - c^2 T(t) X''(x) = 0 \iff T''(x) X(x) = c^2 T(t) X''(x) \iff$$$$\dfrac{T''(t)}{c^2 T(t)} = \dfrac{X''(x)}{X(x)} = - \lm \quad \lm \in \mathbb{R}$$

    Ahora, buscamos resolver:
    \begin{equation*}
        \left\{\begin{array}{l}
            X'' + \lm X = 0 \qquad \text{en\ } [0,L] \\
            X(0) = 0  \\
            X(L) = 0 
        \end{array}\right.
    \end{equation*}
    de donde $X(0)=X(L)=0$ viene de $u(t,0) = 0 = u(t,L)$ para todo $t\in[0,\infty[$, ya que $u(t,0)=T(t)X(0)=T(t)X(L)=u(t,L)=0$ y $T(t)\neq 0$ en general (no para todo valor de $t$) por lo que $X(0)=X(L)=0$ (puede que haya otras condiciones de contorno pero se traducen de la misma forma).\\

    Tenemos entonces un problema de tipo Sturm-Liouville. El polinomio característico asociado es $p(r) = r^2 + \lm$, y sus raíces salen de la ecuación 
    $$p(r) = 0 \iff r^2 + \lm = 0 \iff r^2 = - \lm$$

    Ahora distinguimos, los dos primeros casos, y seguiremos con el restante:
    \begin{enumerate}
        \item Si $\lm < 0$, entonces $r^2 = \lm \iff r = \pm \sqrt{-\lm} = \pm \mu$, para cierto $\mu > 0$.
        La solución general es entonces $X(x) = A e^{\mu x} + B e^{- \mu x}$. Imponiendo $X(0) = 0 = X(L)$:
        \begin{enumerate}
            \item $\boxed{X(0) = 0} \Longrightarrow A + B = 0 \Longrightarrow B = - A$, luego $X(x) = A(e^{\mu x} - e^{- \mu x})$.
            \item $\boxed{X(L) = 0} \Longrightarrow A(e^{\mu L} - e^{- \mu L}) = 0 \stackrel{\mu, L > 0}{\Longrightarrow} A = 0$, luego $X \equiv 0$.
        \end{enumerate}
        \item Si $\lm = 0$, entonces resolvemos $X'' = 0 \Longrightarrow X(x) = mx + n$, para ciertos $m,n \in \mathbb{R}$. Imponemos nuevamente las mismas condiciones de contorno: 
        \begin{enumerate}
            \item $\boxed{X(0) = 0} \Longrightarrow n = 0$, luego $X(x) = mx$.
            \item $\boxed{X(L) = 0} \Longrightarrow mL = 0 \stackrel{L > 0}{\Longrightarrow} m = 0$, luego $X \equiv 0$.
        \end{enumerate}
    \end{enumerate}

    En lo que sigue podemos suponer que $\lm > 0$. En tal caso, la solución general es de la forma 

    \begin{equation*}
        X(x) = c_1\cos(\mu x) + c_2\sen(\mu x), \qquad c_1,c_2\in \mathbb{R}
    \end{equation*}

    Imponemos que $X(0) = 0$, y obtenemos que $c_1 = 0$, luego $X(x) = c_2 \sen (\mu x)$. Ahora, imponemos $X(L) = 0$, de donde
    $$X(L) = c_2 \sen(\mu L) = 0 \stackrel{c_2 > 0}{\Longrightarrow} \sen (\mu L) = 0 \Longrightarrow \mu L = n \pi \stackrel{\mu = \sqrt{\lm}}{\Longrightarrow} \sqrt{\lm} = \dfrac{n \pi}{L} \Longrightarrow \lm_n = \left( \dfrac{n \pi}{L} \right)^2 \quad n \in \mathbb{N}$$
    Entonces, para cada $n \in \mathbb{N}$, tenemos que $X_n(x) = c_2 \sen \left( \dfrac{n \pi x}{L} \right)$. \\

    Pasamos a resolver:
    \begin{equation*}
        \begin{array}{l}
            \dfrac{T_n''(t)}{c^2 T_n(t)} = - \dfrac{n^2 \pi^2}{L^2} \qquad \text{en\ } [0,L]
        \end{array}
    \end{equation*}

    Vemos que:

    $$\dfrac{T_n''(t)}{c^2 T_n(t)} = - \dfrac{n^2 \pi^2}{L^2} \iff T_n''(t) = -  \dfrac{c^2 n^2 \pi^2}{L^2} T_n(t) \iff T_n''(t) + \dfrac{c^2 n^2 \pi^2}{L^2} T_n(t) = 0$$

    El polinomio característico asociado es $q(r) = r^2 + \frac{c^2 n^2 \pi^2}{L^2}$, y sus raíces vienen dadas como las soluciones de

    $$q(r) = 0 \iff r^2 = - \dfrac{c^2 n^2 \pi^2}{L^2} \iff r = \pm i \dfrac{c n \pi}{L}$$

    La solución general es entonces 

    $$T_n(t) = a_n \cos \left( \dfrac{c n \pi t}{L} \right) + b_n \sen \left( \dfrac{c n \pi t}{L} \right) \quad a_n, b_n \in \mathbb{R} \text{ para cada } n \in \mathbb{N}$$

    La solución más general posible, usando que la ecuación es lineal, sabemos que es, llamando $a_n c_1 = A_n$ y $b_n c_1 = B_n$ para cada $n \in \mathbb{N}$, la siguiente:
    $$u(t,x) = \sum_{n=1}^{\infty} u_n(t,x) = \sum_{n=1}^{\infty} \left[A_n \cos \left( \dfrac{c n \pi t}{L} \right) + B_n \sen \left( \dfrac{c n \pi t}{L} \right) \right] \sen \left( \dfrac{n \pi x}{L} \right)$$

    Imponemos condiciones de contorno, y trabajamos con un ejemplo (sencillo) concreto para que se vea más claro:
    
    $$\varphi(x) = \sen \left(\dfrac{2 \pi x}{L} \right) + 2 \sen \left(\dfrac{5 \pi x}{L} \right)$$
    
    de donde
    
    $$u(0,x) = \varphi(x) \iff \sum_{n=1}^{\infty} A_n \sen \left( \dfrac{n \pi x}{L} \right) =  \sen \left(\dfrac{2 \pi x}{L} \right) + 2 \sen \left(\dfrac{5 \pi x}{L} \right) \iff$$$$ A_2 = 1 \quad A_5 = 2 \quad A_n = 0 \quad \forall n \in \mathbb{N} \setminus \{2,5\}$$

    Análogamente, para $\psi(x) = 3 \sen \left( \dfrac{3 \pi x}{L} \right)$ se tiene que:

    $$u_t(t,x) = \sum_{n=1}^{\infty} \left[- \dfrac{c n \pi}{L} A_n \sen \left( \dfrac{c n \pi t}{L} \right) + \dfrac{c n \pi}{L} B_n \cos \left( \dfrac{c n \pi t}{L} \right) \right] \sen \left( \dfrac{n \pi x}{L} \right)$$

    y 

    $$u_t(0,x) = \sum_{n=1}^{\infty} \dfrac{c n \pi}{L} B_n \sen \left( \dfrac{n \pi x}{L} \right) = 3 \sen \left( \dfrac{3 \pi x}{L} \right)$$

    deducimos que, para $n = 3$, $$\dfrac{c 3 \pi}{L} B_3 = 3 \iff B_3 = \dfrac{L}{c \pi}$$

    y además $B_n = 0 \quad \forall n \in \mathbb{N} \setminus \{3\}$. Sustituyendo todo, llegamos a la solución final:
    $$\boxed{u(t,x) = \cos \left( \dfrac{2 c \pi t}{L} \right) \sen \left( \dfrac{2 \pi x}{L} \right) + \dfrac{L}{c \pi} \sen \left( \dfrac{3 c \pi t}{L} \right) \sen \left( \dfrac{3 \pi x}{L} \right) + 2 \cos \left( \dfrac{5 c \pi t}{L} \right) \sen \left( \dfrac{5 \pi x}{L} \right)}$$

    \section*{Anexo. Resolución de Ecuaciones Diferenciales}

    \subsection{Ecuaciones lineales homogéneas}
    Se supone una ecuación de la forma 
    \begin{gather*}
        \sum_{i=0}^n a_i y^{i)} = a_ny^{n)} + ... + a_1 y' + a_0y = 0
    \end{gather*}
    Calculamos el polinomio característico como 
    \begin{gather*}
        \sum_{i=0}^n a_i x^i = a_nx^n + ... + a_1x + a_0
    \end{gather*}
    y extraeremos la raíces de este polinomio. Supongamos que encontramos una raíz $r\in \bb{C}$ que podemos escribir como $r=\alpha + i\beta$ con $\alpha, \beta \in \bb{R}$. De esta forma la solución general de $y$ para esta raíz será:
    \begin{gather*}
        y_r = e^{\alpha x} \left[ C_1 \cos(\beta x) + C_2 \sin(\beta x) \right],\ \ \ C_1,C_2\in \bb{R}
    \end{gather*}
    La solución general $y$ será la combinación lineal de todas estas $y_r$ para cada raíz del polinomio (la suma de ellas valdrá).

    Si obtenemos una raíz real doble $r$, la fórmula general de $y$ (para dicha raíz) será:
    \begin{gather*}
        y_r = C_1 e^{rx} + C_2 x e^{rx}, \ \ \ C_1, C_2\in \bb{R}
    \end{gather*}

    \subsection{Ecuaciones tipo Euler}
    Se supone una ecuación de la forma 
    \begin{gather*}
        \sum_{i=0}^n a_i x^i y^{i)} = a_nx^ny^{n)} + ... + a_1x y' + a_0y = 0
    \end{gather*}
    Suponemos que la solución es de la forma $y(x)=x^r$. De esta forma, tendremos que 
    \begin{gather*}
        y^{i)} = r(r-1)...(r-i+1)x^{r-i} = r_ix^{r-i}\ \ \ \ \forall i \in \{0,...,n\}
    \end{gather*}
    Al sustituir en la ecuación general obtendremos
    \begin{gather*}
        \sum_{i=0}^n a_i x^i y^{i)} = \sum_{i=0}^n a_i x^i r_ix^{r-i} =  \sum_{i=0}^n a_i r_ix^{r} = x^r \sum_{i=0}^n a_i r_i = 0
    \end{gather*}
    Por lo que tendremos que resolver 
    \begin{gather*}
        \sum_{i=0}^n a_i r_i = 0
    \end{gather*}
    ya que $x^r \neq 0$ para cualquier $x\neq 0$. Al resolver esto obtendremos raíces $r\in \bb{C}$ que podemos escribir como $r=\alpha + i\beta$ con $\alpha, \beta \in \bb{R}$. De esta forma la solución general de $y$ para esta raíz será:
    \begin{gather*}
        y_r = x^\alpha [C_1 \cos(\beta \ln x) + C_2 \sin(\beta \ln x)],\ \ \ C_1, C_2\in \bb{R}
    \end{gather*}
    La solución general $y$ será la combinación lineal de todas estas $y_r$ para cada raíz del polinomio (la suma de ellas valdrá).

    Si obtenemos una raíz real doble $r$, la fórmula general de $y$ (para dicha raíz) será:
    \begin{gather*}
        y_r = C_1 x^r + C_2 \ln(x) x^r , \ \ \ C_1, C_2\in \bb{R}
    \end{gather*} 
\end{document}